% history-id: intro
\section{Wprowadzenie}

Rzut ukośny jest przykładem ruchu, w którym prosty model fizyczny prowadzi do
nietrywialnego toru. Ciało otrzymuje prędkość początkową skierowaną pod kątem
do poziomu, a następnie porusza się pod wpływem grawitacji.

W najprostszym modelu pomijamy opór powietrza i traktujemy przyspieszenie
grawitacyjne jako stałe. Celem artykułu jest wyprowadzenie równań ruchu,
równania toru oraz wzoru na zasięg.
